\subsection{Bond prices, yields and spreads}

These entries price a fixed-rate bond from its yield, on and between coupon dates, and turn prices back into yields and spreads, for fixed-rate, floating-rate and money-market instruments. Continuous-compounding zero-rate pricing is in \hyperref[def:interest-rate-measures]{interest-rate measures} and U.S. day counts in \hyperref[conv:us-rate-quotation]{U.S. quotation}. Numbers are checked in \texttt{trading\_notes.numerical.bond\_pricing\_yields.bond\_valuation\_checks}. Prices are per 100 of par and yields bond-equivalent (twice the semiannual rate) unless stated.

\begin{qfmethod}[Pricing a bond from its yield on and between coupon dates: cash flows, the periodic rate, the fractional first period, accrued interest, full and flat prices, and the relations among price, coupon, maturity and yield]\label{meth:bond-price-from-yield}
\qffield{Statement} A bond's price is the present value of its cash flows at the required yield. With semiannual coupons $c/2$ per 100, $n$ coupons remaining and yield $y$, the price on a coupon date is
\[ P=\sum_{k=1}^{n}\frac{c/2}{(1+y/2)^k}+\frac{100}{(1+y/2)^n}. \]
Between coupon dates the street convention discounts each flow over $k-1+w$ periods, where $w$ is the days from settlement to the next coupon over the days in the coupon period, counted by the bond's day-count basis: actual/actual for Treasuries, 30/360 for corporate, agency and municipal bonds. The result is the \emph{full} (dirty) price. The buyer pays it, but the U.S. market quotes the \emph{flat} (clean) price, the full price less accrued interest, $\text{AI}=(c/2)\times(\text{days since last coupon}/\text{days in period})$ \citep[chap.~4, pp.~65 and~72--76]{s7-fabozzifixed}. The Treasury's own formula for a security settling inside a coupon period discounts the stub to the next coupon simply rather than with compounding: $(P+A)[1+(r/s)(i/2)]=C/2+(C/2)a_{\overline n|}+100v^n$, with $r$ days to the next coupon of $s$, $v=1/(1+i/2)$ and $n$ further coupons \citep[sec.~II.D]{cfr-31-356-appendix-b}.

\qffield{Practice} \emph{Reproduced.} The Treasury's examples: an 8\textthreequarters\ percent 30-year bond at 8.84 percent prices at 99.057893, and a 9\textonehalf\ percent note at 9.54 percent two weeks after its dated date at 99.730918 plus accrued interest of 0.367403 \citep[secs.~II.A and~II.D]{cfr-31-356-appendix-b}. Fabozzi's 10 percent corporate bond maturing 1~March~2027, settling 17~July~2021 at 6.5 percent: 44 days of 30/360 to the next coupon, $w=0.24444$, full price 120.0281, accrued 3.7778, flat 116.2503 \citep[chap.~4, Exhibit~4-4, pp.~74--75]{s7-fabozzifixed}.

\emph{Original.} A 5 percent corporate bond (30/360, coupons 15~February and 15~August) maturing 15~February~2034, traded on 20~October~2026 for T+1 settlement on 21~October at a yield of 4.80 percent. From 15~August to 21~October is 66 days of 30/360, so accrued interest is $2.5\times66/180=0.916667$ and $w=114/180=0.633333$; 15 coupons remain. The full price is 102.1316 and the flat price 101.2149; solving back from 102.1316 returns 4.80 percent. Discounting the stub simply, as in the Treasury's formula, gives 102.1249, 0.0067 lower: the convention is part of the quote.

\qffield{Conditions} \emph{Relations} (option-free bonds, on a coupon date). If the coupon equals the yield the bond prices at par; above it at a premium; below it at a discount. At an unchanged yield a premium bond's price falls and a discount bond's rises toward par as maturity approaches (pull to par). Price falls when the yield rises, and the curve is convex: a 100 basis point fall raises the price more than a 100 basis point rise lowers it. Longer maturity and lower coupon raise sensitivity: at 4.80 percent, +100 basis points lowers a 3-year par bond by 2.72 percent, a 15-year par bond by 9.93 percent and a 15-year zero by 13.60 percent (chap.~4, pp.~68--71; chap.~5, pp.~102--105).

\qffield{Pitfalls} Pricing with a yield whose compounding differs from the coupon frequency. Using the wrong day-count basis for accrued interest (30/360 against actual/actual). Comparing flat prices across bonds with very different accrued interest. Bonds in default or income bonds trade flat, without accrued interest (chap.~4, p.~74, n.~5).

\qffield{See also} \hyperref[conv:us-rate-quotation]{U.S. quotation}; \hyperref[def:bond-yield-spread-measures]{yield measures}; \hyperref[meth:yield-duration-pvbp]{duration}.

\end{qfmethod}

\begin{qfmethod}[Matrix pricing: valuing a bond without a reliable quote from the yields or spreads of comparable quoted bonds, the interpolation steps, and its place in the fair-value hierarchy]\label{meth:matrix-pricing}
\qffield{Statement} ``Matrix pricing is a mathematical technique used principally to value some types of financial instruments, such as debt securities, without relying exclusively on quoted prices for the specific securities, but rather relying on the securities' relationship to other benchmark quoted securities''; it is a market-approach valuation technique \citep[IFRS~13, para.~B7]{eu-2023-1803-ifrs13}. An entity holding many similar debt securities may use it as a practical expedient where quoted prices exist but are not readily accessible for each one, but the measurement then falls in a lower level of the fair-value hierarchy (para.~79(a)). Pricing services also use evaluated prices built from issue analysis and dealer quotes \citep[chap.~44, p.~1119]{s7-fabozzifixed}.

\emph{Steps.} (1) Select quoted bonds comparable in issuer type, credit quality, sector, coupon and embedded options. (2) Compute their yields, or their spreads over a benchmark. (3) Interpolate (or extrapolate) to the target bond's maturity, or apply the average spread to the benchmark yield at that maturity. (4) Price the target bond's cash flows at the estimated yield.

\qffield{Practice} Example, original (invented quotes, on a coupon date, semiannual). Bond A, a 4-year 4.00 percent note quoted at 98.20, yields 4.4967 percent; bond B, a 6-year 5.00 percent note quoted at 101.10, yields 4.7869 percent; both share the target's rating and sector. Interpolating linearly gives a 5-year yield of 4.6418 percent, at which an untraded 5-year 4.50 percent bond prices at 99.3737. By the spread route, A and B yield 61.7 and 76.7 basis points over 3.88 and 4.02 percent benchmarks; their mean, 69.2 basis points, over a 3.95 percent 5-year benchmark gives the same 4.6418 percent here, because the benchmark is itself linear between the two maturities.

\qffield{Conditions} The method assumes that yields or spreads vary smoothly with maturity and that the comparables differ from the target only in maturity. Liquidity, size, covenants, call features and issuer-specific credit can all break that assumption; then the output is an estimate with model risk, not a market price.

\qffield{Pitfalls} Interpolating across bonds with different embedded options or seniority. Extrapolating far beyond the comparables' maturities. Treating a matrix price as a Level~1 input or as evidence that the bond could be sold at that price.

\qffield{See also} \hyperref[meth:bond-price-from-yield]{pricing from a yield}; \hyperref[def:bond-yield-spread-measures]{spread measures}; IFRS~13 Level~2 inputs (above).

\end{qfmethod}

\begin{qfdefinition}[Yield and spread measures for fixed-rate bonds: periodicity and compounding conversions, effective annual yield, current yield, yield to maturity, yields to call and to worst, and G-, I- and Z-spreads]\label{def:bond-yield-spread-measures}
\qffield{Statement} \emph{Yields.} The \emph{current yield} is the annual coupon over the price; it ignores capital gain or loss and reinvestment. The \emph{yield to maturity} is the internal rate of return of the full price against the promised cash flows; doubling the semiannual rate gives the bond-equivalent yield. It equals the realised return only if the bond is held to maturity and every coupon is reinvested at that same yield. The \emph{yield to call} uses the cash flows to an assumed call date and call price; the lowest of the yield to maturity and all yields to call (and put) dates is the \emph{yield to worst} \citep[chap.~4, pp.~76--84]{s7-fabozzifixed}. A yield quoted with $m$ periods a year converts to annual compounding as $(1+y_m/m)^m-1$ (pp.~81--82).

\emph{Spreads.} The nominal or G-spread is the yield minus a similar-maturity (often interpolated) government yield (chap.~37, p.~856). The I-spread, called the yield-to-yield asset-swap spread when the benchmark is a matched-maturity par swap rate, is the yield minus that swap rate (chap.~35, p.~814). The Z-spread, also called the static or zero-volatility spread, is the constant spread to every zero rate of a benchmark curve that discounts the cash flows to the full price, $P=\sum_t CF_t/(1+z_t+Z)^t$ (chap.~35, p.~816; chap.~37, pp.~856--857).

\qffield{Practice} \emph{Conversions.} A 6.00 percent semiannual yield is 6.0900 percent annual, 5.9557 percent quarterly, 5.9263 percent monthly and 5.9118 percent continuously compounded.

\emph{The 5 percent 2034 bond} of \hyperref[meth:bond-price-from-yield]{pricing from a yield} (full price 102.1316, yield 4.80 percent): current yield 4.94 percent. With a call schedule of 102 in February~2029, 101 in February~2030 and par from February~2031, the yields to call are 5.2558, 4.8779 and 4.6841 percent; the yield to worst is the 2031 par call, 4.6841 percent.

\emph{Spreads} over an invented government zero curve $z(t)=3.60\%+0.08\%\,t$ (semiannual compounding), which gives 7- and 10-year par yields of 4.1333 and 4.3425 percent: G-spread 64.4 basis points over the interpolated 7.3-year par yield; Z-spread 65.0 basis points; and against SOFR swap rates of 3.80 and 3.95 percent, an I-spread of 98.4 basis points. For a near-par bullet bond the G- and Z-spreads stay within a basis point even on a steeper curve ($2.50\%+0.30\%\,t$). For a 7.3-year zero-coupon bond at the same 4.80 percent yield on that steeper curve they differ: G-spread 23.0 basis points, Z-spread 10.4, because the zero's single cash flow is discounted at the 4.70 percent zero rate, not at a par yield.

\qffield{Conditions} Yield to maturity assumes reinvestment at the yield and holding to maturity; for an 18-year 6 percent bond bought at 9.5 percent, interest on interest supplies about 54 percent of the total dollar return the yield promises (p.~79). Yield to call assumes the call is exercised on that date, which a callable bond's value does not; option-adjusted measures need a rate model (p.~84).

\qffield{Pitfalls} Doubling a semiannual rate and calling it an annual effective yield (bond-equivalent understates effective). Averaging the yields of a portfolio's bonds instead of computing its cash-flow yield (\hyperref[meth:convexity-portfolio-duration]{portfolio measures}). Quoting a G-spread for amortising cash flows; the Z-spread uses the whole curve.

\qffield{See also} \hyperref[meth:matrix-pricing]{matrix pricing}; \hyperref[def:swap-curve-short-term-spreads]{swap spreads}; \hyperref[meth:bond-return-decomposition]{horizon return}.

\end{qfdefinition}

\begin{qfmethod}[Floating-rate notes and money-market yields: quoted and required margins, discount margin and other floater spreads, the Treasury FRN formula; discount-basis, money-market (add-on) and bond-equivalent (investment-rate) yields and their conversions]\label{meth:floating-money-market-yields}
\qffield{Statement} \emph{Floaters.} A floater pays a reference rate plus a \emph{quoted margin} fixed at issue; the \emph{required margin} is the spread the market now demands. On a reset date the floater prices at par if the two are equal, below par if the required margin is higher. Its price also moves with the time to the next reset and with any cap or floor \citep[chap.~14, pp.~312--318]{s7-fabozzifixed}. Spread measures include spread for life, adjusted simple and total margins, and the \emph{discount margin}, the spread over the current reference rate that, with the reference rate assumed unchanged, discounts the projected cash flows to the price (pp.~315--316; chap.~4, pp.~85--86). The Treasury prices its FRNs this way, accruing daily at $\max(r+s,0)/360$ on the index rate $r$ and spread $s$, and discounting each quarterly flow at $1+(\text{days}/360)(r+m)$ with discount margin $m$ \citep[sec.~IV.A]{cfr-31-356-appendix-b}.

\emph{Money-market yields.} Treasury bills and commercial paper are quoted on a \emph{discount basis} against face value over a 360-day year: price $P=100(1-d\,t/360)$ for $t$ days \citep[sec.~VI.A]{cfr-31-356-appendix-b}; \citep[chap.~13, p.~301]{s7-fabozzifixed}. A \emph{money-market} (add-on, CD-equivalent) yield is $(100-P)/P\times360/t$; the Treasury's \emph{investment rate} (bond-equivalent yield) is $(100-P)/P\times y/t$ with $y$ 365 or 366 days for bills of up to half a year, and solves a quadratic for longer bills \citep[sec.~VI.D]{cfr-31-356-appendix-b}.

\qffield{Practice} \emph{Reproduced.} Fabozzi's six-year floater paying the reference rate plus 80 basis points semiannually, reference rate 10 percent, priced 99.3098: discount margin 96 basis points (48 per half-year) \citep[chap.~4, Exhibit~4-7, p.~86]{s7-fabozzifixed}. The Treasury's bills: a 7.610 percent discount for 90 days is a price of 98.0975; a 182-day bill at 95.934567 has a discount rate of 8.042 percent; a 20-day bill at 99.559444 an investment rate of 8.076 percent; a 364-day bill at 92.265 an investment rate of 8.237 percent by the quadratic \citep[secs.~VI.A, C and~D]{cfr-31-356-appendix-b}. Its FRN index is the 13-week bill auction rate restated as a simple 360-day yield: a 0.095 percent discount over 91 days becomes 0.095022819 percent (sec.~IV.A).

\emph{Original.} A 182-day bill at a 4.00 percent discount prices at 97.977778: money-market yield 4.0826 percent, investment rate 4.1393 percent, effective annual 4.1822 percent; for short bills the ordering is discount rate below money-market yield below investment rate. A 4.20 percent act/360 add-on rate is 4.2583 percent on a 365-day basis.

\emph{Original SOFR FRN.} A three-year note issued and priced on 21 October 2026 pays quarterly on the 21st of January, April, July and October (modified following), 12 periods of 90 to 94 days. Its coupon is SOFR compounded in arrears over each period, with no lookback, plus a 75 basis point margin, act/360. Compounding follows the New York Fed's SOFR Averages arithmetic: one factor $1+\text{SOFR}_in_i/360$ per business day, $n_i$ the calendar days until the next business day, so a Friday rate counts for three days in a single simple-interest term \citep[``Calculation Methodology for the SOFR Averages and Index'']{nyfed-sofr-averages-index}. With SOFR flat at 4.10 percent, the first 92-day period grows by 1.010532, a compounded rate of 4.1211 percent, and pays 1.053160 plus 0.191667 margin, 1.244827 per 100.

Each period is discounted at its projected SOFR growth plus a discount margin $m$ accrued simply, $1/(G_k+m\,d_k/360)$, so the note is at par on a reset date when $m$ equals the quoted margin. At 99.40 it has a discount margin of 96.38 basis points: it trades below par because the market requires more than its quoted margin. Equivalently, $P=100+100(0.75\%-m)A$, where $A=\sum_k(d_k/360)D_k=2.8062$ is the annuity of the discount factors.

\qffield{Conditions} The discount margin assumes the reference rate stays at its current level and ignores caps, floors and calls (chap.~4, p.~87; chap.~14, p.~316). A compounded SOFR coupon is known only at the end of its period, unlike a LIBOR fixing known at the start (\hyperref[def:interest-rate-swaps]{SOFR conventions}); issued notes usually observe SOFR with a lookback or observation shift of two to five business days so that the coupon is known before payment. The example's business days are weekdays other than Federal Reserve holidays; SOFR itself is not published for days on which SIFMA recommends a full closure of U.S. government securities trading, a calendar that can differ \citep[``Details on Publication and Revisions for the Treasury Repo Reference Rates'']{nyfed-sofr-averages-index}. Bill quotations, day bases and year lengths are U.S. conventions; other markets quote money-market yields on 365-day bases.

\qffield{Pitfalls} Comparing a bill's discount rate directly with a coupon bond's yield; it is computed on face value and a 360-day year, so it understates the return. Treating a floater as free of price risk: margin changes, reset lags and caps move its price.

\qffield{See also} \hyperref[conv:us-rate-quotation]{U.S. quotation}; \hyperref[def:money-market-reference-rates]{money-market reference rates}; \hyperref[def:bond-yield-spread-measures]{fixed-rate yields}.

\end{qfmethod}
